\chapter*{Abstract}
\addcontentsline{toc}{chapter}{Abstract}

Complete-outfit recommendation requires a system to model relationships among
garments while respecting practical constraints and the item already selected
by a user. This study evaluates whether a small project-specific compatibility
head can support that task without retraining a large visual encoder or tying
the output to a single retailer catalogue. Frozen 768-dimensional SigLIP image
representations were combined with a 379,265-parameter compatibility head and
trained on the published Polyvore Outfits disjoint compatibility split. Here,
\emph{lightweight} refers to the trainable head and its downstream training
cost, not to the size or original pretraining cost of SigLIP. The
model was compared with non-learned, linear, recurrent, type-aware and
attention-based controlled baselines over three random seeds. Evaluation included
classification, corrected FITB ranking, calibration, structural ablation,
imbalance, overfitting, runtime, bias and domain shift.

The selected compatibility model achieved a held-out ROC--AUC of 0.8583 (95\%
bootstrap confidence interval 0.8524--0.8641) and 60.25\% FITB accuracy (95\%
confidence interval 58.73--61.75). The three-seed mean FITB result was 59.94\%.
Removing the pairwise relationship module reduced mean FITB by 10.28 percentage
points, with a paired 95\% interval of 8.96--11.57 points; the smaller gain from
slot embeddings was inconclusive. Training AUC was 0.9699. After
validation-based early stopping, validation and test AUC differed by only
0.00145.

The model was integrated into an abstract, product-independent candidate space
with weather and clothing-range constraints. Exact search exposed substantial
loss from the previous quota path, which recovered the exact top result in
27.78\% of 18 controlled queries. The findings support the feasibility of an
auditable system with lightweight downstream training within the tested
setting. Evaluation was limited by 29.50\% official FITB coverage, sparse
pure-menswear evidence, residual item-identifier overlap across the published
split files, strong catalogue domain shift and the absence of participant
evaluation.

\textbf{Keywords:} fashion recommendation; outfit compatibility; SigLIP;
fill-in-the-blank; constrained search; responsible artificial intelligence

\clearpage
